\section{評価と考察}

ここに各評価仕様について記す.
評価は伝達精度と伝達速度,学習用としてわかりやすいかどうかで評価を行う.


まず伝達精度は送信ビットと受信ビットを比較し,1文字あたりのビット誤り率を精度としている.
音のシステムの精度目標はFFTのポイント数を512に設定した時の周波数検出率より90 \%を目標としている\cite{otoken}.
光のシステムの精度目標はLEDと照度センサの距離が10 cm未満のため100 \%を目標としている\cite{hitensou}.
球のシステムの精度目標はサーボモータの静止角度誤差が±5 \%のため1つのサーボモータの分岐成功率を95 \%と見積もる\cite{seisi}.
使用するサーボモータの数は8つのため$(0.95)^8 = 0.663$より66 \%を目標としている.

次に伝達速度は送信開始から受信終了まで1文字あたり何秒かかるかを速度としている.
伝達速度の目標は人間のシステムの応答を待てる時間が2000 msのため音,光,球それぞれ2000 msを目標とする\cite{outou}.

最後にわかりやすいかどうかはモールス信号の書き取り精度と各システムの書き取り精度を比較する.
各システムの書き取り精度がモールス信号の書き取り精度を上回ることを目標とする.


結果より音の伝達精度は85 \%,伝達速度は平均2476.7 ms,書き取り精度は86 \%だった.
光の伝達精度は100 \%,伝達速度は平均817.91 ms,書き取り精度は16 \%だった.
球の伝達精度は100 \%,伝達速度は平均4172.0 ms,書き取り精度は90 \%だった.
またモールス信号の書き取り精度は42 \%だった.
よって評価目標は音の伝達精度,伝達速度,光の書き取り精度,球の伝達速度以外は目標を上回る事ができた.
また球のシステムのセンサが動かなかった.


まず音の評価を行う.
音は伝達精度と伝達速度が評価指標を上回らなかった.
音の伝達精度は指標の90\%に対して平均85\%と5\%指標を下回った.
この原因としては試行をおこなった場所が環境ノイズのある場所であったためだと考えられる.
今回使用した周波数は440Hzから587Hzであるがそれは男性の話し声の高さである500Hzと近しいため,誤判別が多かった\cite{koe}.
よって人間の話し声でなく且つ可聴できる周波数の2500Hzなどの本システムよりも高い周波数にするのが良いと考えられる.
そして音の伝達速度は指標の2000 msに対して476.7 ms指標を超えてしまった.
この原因としてはシステムのパラメータとして音の長さを500 msと設定しており,4音で1文字を表しシステムの動く時間を含めずに2000 msに達していたためだと考えられる.
よって速度を上げるためには音の長さを500 msから400 msに変更しシステムの動く時間や待機時間の余裕を作るべきだ.


次に光の評価を行う.
光は書き取り精度が評価目標を上回らなかった.
モールス信号の書き取り精度が42 \%に対して光の書き取り精度は16 \%であった.
この原因としては点灯の間に消灯する空白時間がなかったため一度光ったものが何個分の1を表しているかがとても分かりづらかったことが考えられる.
よってこの対策として点灯の間に短い消灯する時間を設ける必要がある.
さらに連続する消灯の個数も分かりづらかったと考えられるため,信号の間を消灯,0というビットを送りたい場合は今のシステムの点灯の半分の明るさの点灯に変更するのがよいと考えられる.

%%その理由としては自然光に左右されないように動的に閾値を補正するプログラムを使用することにより周りの環境によらずに行えたためだと考えられる.
また段ボールで覆ってしまったことで段ボールの隙間からしか伝達の様子が可視できなくなってしまったことは課題である.
よって段ボールにシステムに支障のない程度の小さな穴を開けておいてそこから伝達の様子を見るというのが良いと考えられる.


最後に球の評価を行う.
球は伝達速度が評価指標を上回らなく,受信の照度センサが動かなかった.
球の伝達速度は指標の2秒に対して平均4.1秒と2倍以上の遅さであった.
この原因としては指標が直線距離40 cmの等加速度運動で計算したのに対して,実際は分岐点で速度が減少してしまったり最短距離やそれに近いルートではない回り道をしていた点が挙げられる.
球は必ず3回分岐点を通らなくてはいけない.
その度に減速しているため遅くなってしまったと考えられる.
距離は加速度の2乗に比例して大きくなるため,大幅に遅れてしまっている.
また照度センサが動かなかった.
その原因としては球がセンサの上部を通ったときに変化する電圧が通っていない時のノイズ電圧の変化よりも小さいことである.
試行の後指で横から球を支えて穴の上で停止させたときはセンサが機能していたことから,球が完全に穴を覆うことができればセンサが反応し数字を返すことができるが,伝達速度向上のため板の角度を急にしたことによって球の通過速度が早すぎて完全に覆いきれなかったと考える.
また球の通る道が球の大きさよりも必要以上に大きく穴の真上を通らなかったことも原因に挙げられる.
よって穴付近の球の通る道を球が突っかからない程度の広さにし,板の角度を緩やかにすることでうまく機能させるのが良いと考えられる.